\section{Discussion}\label{disc}
\par % summary of findings
We sought to clarify several concepts relevant to
characterizing the distribution of durations selling or buying sex.
First, we distinguished between "source" and "active" populations,
with the active population being a length-biased sample of the source population.
Second, we reiterated the familiar implications of randomly censored durations,
which halves the mean duration.
We combined the above to derive the exact relationship between
the distribution of censored durations among the active population and
the distribution of total durations among the source population,
which happen to be exactly equal if duration are Exponentially distributed.
We then developed a parametric Bayesian model of this relationship
to infer the distribution of total durations among the source population
from aggregate data on censored durations among the active population.
Applying this model to data from a previous global systematic review,
we inferred significantly shorter mean durations across multiple regions.
\par % implications for HIV incidence
Our estimates of mean duration selling sex were
around half or less versus those reported in \cite{Fazito2012}.
This was expected, as we showed how both methods 2 and 3 can overestimate mean duration,
especially if surveys undersample individuals in the early years of sex work.
Shorter durations suggest that HIV incidence in sex work
could be higher than previously estimated via transmission models,
since there is less time for infections to accumulate
to match observed HIV prevalence in sex work \cite{Knight2020}.
However, existing model-based estimates of HIV incidence
among female sex workers in Sub-Saharan Africa
already exceed empirical estimates \cite{Stevens2024inc};
this might be partly explained by recruitment of
relatively lower-risk female sex workers and
relatively higher-risk wider population women
for empirical HIV incidence studies \cite{Jones2024}.
% JK: not sure where to go with this ...
Rapid sex work turnover could also suggest that
more infections in the wider population might be attributable to
unmet prevention needs during brief periods of sex work \cite{Knight2020}.
\par % implications for modelling
We also inferred that distributions of total durations in sex work are likely "heavy-tailed",
such that rates of exit decline with duration in sex work.
This declining exit rate could reflect
selection effects acting on between-person differences in propensity to remain in sex work,
\ie "frailties" \cite{Balan2020}, and/or
within-person changes while in sex work,
including self-identification as a sex worker, and
changing social ties with family, friends, and clients
\cite{Neufeld2023,Bhattacharjee2023,Ingabire2012,Bowen2011,Lora2025}.
While individual-based simulation models could
specify "heavy-tailed" distributions directly,
compartmental simulation models would need
two or more sex work states to reproduce this pattern --- \eg
separate short and long duration sex work trajectories, and/or
separate early and late sex work periods.
Indeed, such states would also be useful to capture
unique risks associated with early sex work
\cite{Neufeld2023,Bhattacharjee2023}.
% TODO: revise above to note that models can/should model particular determinants?
% ------------------------------------------------------------------------------
\paragraph{Limitations}
Our work has several limitations.
First, we did not model any gaps (temporary breaks) in sex work.
However, many estimates of sex work duration are derived from
the difference between respondents' current age and the age they report first selling sex,
which implicitly ignores such gaps.
Respondents may also report "duration selling sex" in this way when asked directly,
especially since unsuccessful attempts to exit sex work can be stigmatized \cite{Lora2025}.
% JK: Should we recommend that modellers try to incorporate these "hidden gaps" somehow?
%     Frankly I'm not sure what specifically I would recommend.
Second, we did not consider any recall/reporting biases,
as entry into sex work is likely a memorable life-course event
\cite{Sandy2025,Cimino2017,Ingabire2012,Bowen2011,Lora2025},
and self-identified sex workers do not, to our knowledge, have
obvious motivations for under-reporting durations.
We additionally assumed that sex work turnover is stable over time,
despite empirical evidence of changing sex work dynamics in several contexts
\cite{Anderegg2024,Anderson2026}.
Finally, our re-analysis of duration data from \cite{Fazito2012}
did not include more recent or individual-level data,
since our motivation was mainly to compare
results using our proposed versus prior estimation methods.
Additionally, in forthcoming work we will synthesize
individual-level data from 183 studies \cite{Anderson2026}.
% ------------------------------------------------------------------------------
\paragraph{Future work}
Our Bayesian model could be extended in several ways.
First, it might be more useful to model sex work durations
via semi/parametric hazard rates of exiting sex work,
rather than parametric duration distributions.%
\footnote{Indeed, the hazard rates implied by
  some parametric distributions we considered may be implausible,
  \eg Log-Normal.}
A hazard rate approach could allow and incorporation of
more intuitive parameter priors (\eg 0.5--50\% exit rate per month),
covariates (\eg under proportional hazards assumptions) \cite{Kalbfleisch2023}, and
frailties (\eg to capture unmodelled heterogeneity) \cite{Balan2020}.
Relatedly, joint modelling of duration and age and/or calendar time
could help characterize relevant trends noted above
\eg for HIV/STI transmission modelling \cite{Anderegg2024,Anderson2026}.
Such extensions might be supported by related work estimating
durations (also called sojourn times) in chronic disease states
for modelling and analysis of screening programs \cite{Zelen1969}.
% ------------------------------------------------------------------------------
\paragraph{Conclusion}
